% Revize 2026-10-27: cílené doplnění; nové úlohy a řešení jsou vlastní tvorba.
% Metodický vzor: 4_rocnik/pripravy/Uvod_do_kombinatoriky.pdf (samostatný dokument).
% Hodina 33, 1. ročník SPŠ. Rozsah: Rozpis_vyuky_MAT_1_rocnik.tex, řádek 33.

% Plán: 1_rocnik/tematicke_plany/SPS.pdf; společná projektová šablona.

% Učebnice: ucebnice/01_Zakladni_poznatky/Mocniny_a_odmocniny/Mocniny a Odmocniny.tex

% Původ úloh: Vlastní zadání a vlastní formulace výkladu; metodická inspirace zdroji.

% Metodický zdroj: externi_zdroje/realisticky/Odmocniny/07_Druha_odmocnina.pdf

% URL: http://www.realisticky.cz/ucebnice/01%20Matematika%20S%C5%A0/01%20Z%C3%A1kladn%C3%AD%20poznatky/02%20%C4%8C%C3%ADseln%C3%A9%20obory/07%20Druh%C3%A1%20odmocnina.pdf

% Metodický zdroj: externi_zdroje/realisticky/Odmocniny/08_Treti_odmocnina.pdf

% URL: http://www.realisticky.cz/ucebnice/01%20Matematika%20S%C5%A0/01%20Z%C3%A1kladn%C3%AD%20poznatky/02%20%C4%8C%C3%ADseln%C3%A9%20obory/08%20T%C5%99et%C3%AD%20odmocnina.pdf

\documentclass[11pt,a4paper]{article}

% Spolecny styl pro ucitelske pripravy MAT (Work/MAT)
% Verze 1.1, 2026-10-23
\NeedsTeXFormat{LaTeX2e}
\ProvidesPackage{priprava}[2026/10/23 v1.1 Sablona ucitelskych priprav MAT]

% Zakladni typografie a cestina
\RequirePackage[T1]{fontenc}
\RequirePackage[utf8]{inputenc}
\RequirePackage[czech]{babel}
\RequirePackage{lmodern}
\RequirePackage{microtype}
\RequirePackage{amsmath,amssymb}
\RequirePackage{xcolor}
\RequirePackage{enumitem}
\RequirePackage[most]{tcolorbox}
\RequirePackage{tabularx,booktabs,array}
\RequirePackage{fancyhdr}
\RequirePackage{needspace}
\RequirePackage{geometry}
\RequirePackage[hidelinks]{hyperref}

\geometry{a4paper,top=17mm,bottom=18mm,left=23mm,right=23mm}
\definecolor{notegray}{gray}{0.39}
\setlength{\parindent}{0pt}
\setlength{\parskip}{2.5pt plus 1pt minus 0.5pt}
\setlist[itemize]{topsep=2pt,itemsep=1pt,parsep=0pt,leftmargin=1.45em}
\setlist[enumerate]{topsep=2pt,itemsep=2pt,parsep=0pt,leftmargin=1.65em}
\setcounter{secnumdepth}{0}
\renewcommand{\arraystretch}{1.12}

% Zahlavi: vlevo rocnik | globalni tema, vpravo cislo: nazev hodiny
% Zapati zustava zcela prazdne.
\pagestyle{fancy}
\fancyhf{}
\renewcommand{\headrulewidth}{0.35pt}
\renewcommand{\footrulewidth}{0pt}
\setlength{\headheight}{14pt}
\newcommand{\nastavPripravu}[4]{%
  \fancyhead[L]{\small #1\enspace\textbar\enspace #2}%
  \fancyhead[R]{\small #3: #4}%
}

% Minimalisticke cernobile boxy. Kazda definice zvlast.
\tcbset{commonlesson/.style={enhanced,breakable=false,colback=white,colframe=black,
  colbacktitle=white,coltitle=black,boxrule=0.45pt,
  arc=0.8mm,outer arc=0.8mm,left=3mm,right=3mm,top=1.6mm,bottom=1.6mm,
  toptitle=1mm,bottomtitle=0.7mm,fonttitle=\bfseries,
  attach boxed title to top left={xshift=2.5mm,yshift=-1.15mm},
  boxed title style={colback=white,colframe=white,boxrule=0pt,left=1mm,right=1mm}}}
\newtcolorbox{zadaniBox}[1][]{commonlesson,title={Zadání},#1}
\newtcolorbox{prikladBox}[1][]{commonlesson,title={Příklad},#1}
\newtcolorbox{definiceBox}[1][]{commonlesson,title={Definice},#1}

% Spolecny box pro zadani a jeho reseni, bez mezery mezi castmi.
% Pouziti:
%   \begin{prikladReseniBox}
%     Zadani ulohy.
%     \tcblower
%     \begin{align*} ... \end{align*}
%   \end{prikladReseniBox}
% Spodni cast ma jemny sedy podklad a delici linku; lze v ni sazet
% vice radku matematiky. Obe casti maji jen jeden spolecny ramecek.
\tcbset{solutionlesson/.style={
  commonlesson,
  enhanced jigsaw,
  breakable,
  colbacklower=black!3,
  segmentation style={solid,draw=black!25,line width=0.35pt},
  fontlower=\small\color{notegray},
  before lower={\noindent\textbf{Řešení:}\enspace},
  leftlower=3mm,rightlower=3mm
}}
\newtcolorbox{prikladReseniBox}[1][]{solutionlesson,title={Příklad},#1}
\newtcolorbox{zadaniReseniBox}[1][]{solutionlesson,title={Zadání},#1}

% Pomucky pro autora pripravy: sedive poznamky; \reseni pro starsi pripravy
\newcommand{\poznamka}[1]{\par\begingroup\small\color{notegray}\itshape
  \textnormal{\textbf{Poznámka pro učitele.}}\enspace #1\par\endgroup}
\newcommand{\reseni}[1]{\par\begingroup\small\color{notegray}
  \textbf{Očekávané řešení:}\enspace #1\par\endgroup}
\newcommand{\blok}[1]{\needspace{5\baselineskip}\subsection*{#1}}
\newcommand{\vyrok}[1]{\uv{#1}}


\nastavPripravu{1. ročník}{Odmocniny}{33}{Pravidla pro odmocniny}

\begin{document}

\blok{Výklad a souvislosti}

Odmocninu součinu můžeme zkoumat pomocí definice: součin dvou nezáporných odmocnin je nezáporný a jeho druhá mocnina je původní součin. Podmínky jsou součástí tvrzení, nikoli dodatkem k výsledku.

\begin{definiceBox}[title={Připomenutí definice},breakable=false]
\textbf{Druhá odmocnina} $\sqrt a$ je pro $a\ge0$ nezáporné číslo, jehož druhá mocnina je $a$.
\end{definiceBox}


Pro $a,b\ge0$ platí $\sqrt{ab}=\sqrt a\sqrt b$ a pro $a\ge0$, $b>0$ platí $\sqrt{a/b}=\sqrt a/\sqrt b$. Rovněž $\sqrt{a^2}=|a|$ pro všechna reálná $a$. Pro třetí odmocniny platí odpovídající pravidla pro všechna reálná $a,b$, u podílu s $b\ne0$.

Odvození součinu: $(\sqrt a\sqrt b)^2=ab$ a $\sqrt a\sqrt b\ge0$, takže definice určuje právě $\sqrt{ab}$. Pro součet takové pravidlo neplatí.

\blok{Řešené příklady}

\begin{prikladReseniBox}[title={Příklad 1},breakable=false]
Vypočítej $\sqrt{36\cdot25}$ a $\sqrt[3]{8/27}$.
\tcblower
$\sqrt{36}\sqrt{25}=6\cdot5=30$; $\sqrt[3]8/\sqrt[3]{27}=2/3$. Druhou hodnotu ověříme umocněním na třetí.
\end{prikladReseniBox}


\begin{prikladReseniBox}[title={Příklad 2},breakable=false]
Uprav $\sqrt{x^2}$ pro reálné $x$ a rozhodni, zda $\sqrt{9+16}=\sqrt9+\sqrt{16}$.
\tcblower
$\sqrt{x^2}=|x|$; například pro $x=-3$ je výsledek $3$. U součtu je levá strana $5$, pravá $7$, rovnost tedy obecně neplatí.
\end{prikladReseniBox}


\blok{Samostatné procvičení}

\begin{zadaniReseniBox}[title={Procvičení},breakable=false]
\begin{enumerate}[label=\alph*)]
\item Vypočítej $\sqrt{12}\sqrt3$.
\item Vypočítej $\sqrt{75/3}$ a $\sqrt[3]{(-16)/2}$.
\item Pro jaká reálná $x$ lze psát $\frac{\sqrt x}{\sqrt{x+1}}=\sqrt{\frac{x}{x+1}}$?
\end{enumerate}
\tcblower
\begin{enumerate}[label=\alph*)]
\item $\sqrt{36}=6$; oba radikandy jsou nezáporné.
\item $\sqrt{25}=5$ a $\sqrt[3]{-8}=-2$.
\item Levá strana vyžaduje $x\ge0$ a $x+1>0$, dohromady $x\ge0$. Na tomto oboru platí uvedené pravidlo. Pravá strana sama má i další přípustné hodnoty, které však nepřidáváme k oboru původního výrazu.
\end{enumerate}
\end{zadaniReseniBox}


\poznamka{Do obecného pravidla druhé odmocniny nedosazuj záporné činitele, i když je jejich součin kladný: $\sqrt{(-2)(-8)}=4$, ale $\sqrt{-2}\sqrt{-8}$ v reálných číslech smysl nemá. Zkouška znaménkem je nezbytná.}


\blok{Navazující procvičování ve třídě}

\begin{zadaniReseniBox}[title={Základ: podmínky jsou součástí pravidla},breakable=false]
\begin{enumerate}[label=\alph*)]
\item Vypočítej $\sqrt8\sqrt{18}$, $\frac{\sqrt{147}}{\sqrt3}$ a $\sqrt[3]{-4}\sqrt[3]{16}$.
\item Je správný krok $\sqrt{(-2)(-8)}=\sqrt{-2}\sqrt{-8}$ v reálných číslech?
\item Urči podmínky $\frac{\sqrt{x+2}}{\sqrt x}$ a sluč odmocniny.
\end{enumerate}
\tcblower
\begin{enumerate}[label=\alph*)]
\item $\sqrt{144}=12$, $\sqrt{49}=7$, $\sqrt[3]{-64}=-4$.
\item Ne. Levá strana existuje a je $4$, pravá v $\mathbb R$ neexistuje. Pro rozdělení druhé odmocniny musí být oba činitelé nezáporní.
\item $x>0$; pak $\sqrt{\frac{x+2}{x}}$. Původní přísnější podmínku zachováme.
\end{enumerate}
\end{zadaniReseniBox}

\begin{zadaniReseniBox}[title={Rezerva: sudá a lichá odmocnina},breakable=false]
\begin{enumerate}[label=\alph*)]
\item Uprav $\sqrt{a^2b^2}$ a $\sqrt[3]{a^3b^3}$ pro reálná $a,b$.
\item Vypočítej $\frac{\sqrt{(-3)^2}}{\sqrt{9}}$ a vysvětli, proč výsledek není $-1$.
\end{enumerate}
\tcblower
\begin{enumerate}[label=\alph*)]
\item $|ab|=|a||b|$ a $ab$. Druhá odmocnina je nezáporná, třetí zachovává znaménko.
\item $3/3=1$; ve jmenovateli i čitateli je nezáporná odmocnina.
\end{enumerate}
\end{zadaniReseniBox}


\poznamka{Základ: původní pravidla a první nová skupina od čísel k proměnné. Druhá skupina je rezerva; u absolutní hodnoty vždy vyzkoušej i záporný činitel. Domácí pětice obsahuje kontrolu podmínek i protipříklad. Částečné odmocňování zatím soustavně netrénujeme, to je náplň příští hodiny.}

\blok{Domácí procvičení}

\begin{zadaniReseniBox}[title={Domácí úlohy},breakable=false]
\begin{enumerate}[label=\alph*)]
\item Vypočítej $\sqrt{15}\sqrt{60}$ a $\frac{\sqrt{180}}{\sqrt5}$.
\item Vypočítej $\sqrt[3]{-12}\sqrt[3]{18}$.
\item Uprav $\sqrt{(x-2)^2}$.
\item Urči podmínky a sluč $\frac{\sqrt{x-1}}{\sqrt{x+1}}$.
\item Vyvrat rovnost $\sqrt{a+b}=\sqrt a+\sqrt b$ konkrétní dvojicí nezáporných čísel.
\end{enumerate}
\tcblower
\begin{enumerate}[label=\alph*)]
\item $30$ a $6$.
\item $\sqrt[3]{-216}=-6$.
\item $|x-2|$ pro všechna reálná $x$.
\item $x\ge1$, výsledek $\sqrt{\frac{x-1}{x+1}}$ s touto podmínkou.
\item Například $a=9,b=16$: $5\ne7$.
\end{enumerate}
\end{zadaniReseniBox}

\end{document}
